\documentclass{notaaula}

        \titulonota{Forças, gravitação e órbitas}
        \creditos{Turma 1B · Física · Setembro/2026 · Prof. Josué Cavalcante}

        \begin{document}
        \cabecalho

        \begin{sobre}
        Dinâmica: força peso, força normal, atrito, gravitação universal e leis de Kepler. Habilidades em foco: EM13CNT207, EM13CNT301, EM13CNT306 e EM13CNT307.
        \end{sobre}

        \section{Forças de contato e força peso}

\begin{copiar}{Diagrama de forças}
\definicao Uma força representa uma interação. Antes de calcular, isole o corpo e desenhe
apenas as forças que atuam nele. A resultante determina a aceleração:
\[
  \vec F_{\mathrm{res}}=m\vec a.
\]
A força peso é a atração da Terra sobre o corpo:
\[
  \vec P=m\vec g.
\]
Perto da superfície, adotaremos $g=10\un{m/s^2}$ quando o enunciado não informar outro valor.
A força normal $\vec N$ é perpendicular à superfície; ela não é automaticamente igual ao peso.
\simbolos{$F$, $P$ e $N$ (N); $m$ (kg); $a$ e $g$ ($\mathrm{m/s^2}$).}
\end{copiar}

\begin{copiar}{Atrito}
O atrito se opõe ao deslizamento ou à tendência de deslizamento entre superfícies.
\[
  f_{\mathrm{est}}\leq \mu_e N
  \qquad
  f_{\mathrm{cin}}=\mu_c N.
\]
O atrito estático se ajusta até um valor máximo; o cinético atua quando já há escorregamento.
\atencao{Não escreva sempre $f=\mu N$. No repouso, o atrito estático pode ser menor que $\mu_eN$.}
\diaadia{Pneus, solas e freios dependem de atrito. Água ou óleo podem reduzir a aderência.}
\end{copiar}

\begin{exemplo}
Um bloco de $10\un{kg}$ desliza em piso horizontal com $\mu_c=\dec{0,20}$.
Uma força horizontal de $30\un{N}$ atua para a direita. Determine a aceleração.

\resolucao Como não há aceleração vertical, $N=P=mg=100\un{N}$.
Logo, $f_{\mathrm{cin}}=\mu_cN=20\un{N}$. A resultante horizontal vale
$30-20=10\un{N}$ e
\[
  a=\frac{F_{\mathrm{res}}}{m}=\frac{10}{10}
  =\resultado{1\un{m/s^2}}.
\]
\end{exemplo}

\section{Gravitação universal}

\begin{copiar}{Atração entre massas}
Duas massas se atraem com forças de mesmo módulo e sentidos opostos:
\[
  F=G\frac{m_1m_2}{r^2},
  \qquad G=\dec{6,67}\times10^{-11}\un{N\,m^2/kg^2}.
\]
A distância $r$ é medida entre os centros dos corpos. Para um astro de massa $M$ e raio $R$,
\[
  g=G\frac{M}{R^2}.
\]
Se a distância dobra, a força cai para $1/4$; se triplica, cai para $1/9$.
\end{copiar}

\begin{exemplo}[Exemplo de proporção]
Um satélite sofre força gravitacional $F$ a uma distância $r$ do centro da Terra.
A $2r$, mantendo as massas, a nova força é
\[
  F'=G\frac{Mm}{(2r)^2}=\resultado{\frac{F}{4}}.
\]
\end{exemplo}

\section{Leis de Kepler}

\begin{copiar}{Órbitas planetárias}
\begin{itens}
\item \dest{1ª lei:} a órbita é elíptica, com o Sol em um dos focos.
\item \dest{2ª lei:} a linha Sol--planeta varre áreas iguais em tempos iguais. O planeta se move
      mais rápido quando está mais próximo do Sol.
\item \dest{3ª lei:} para corpos que orbitam o mesmo astro,
      $T^2/a^3$ é constante, em que $a$ é o semieixo maior.
\end{itens}
Essas leis descrevem o movimento; a gravitação de Newton explica a interação que mantém a órbita.
\end{copiar}

\begin{dica}
Em queda livre, a ausência de apoio produz sensação de falta de peso. A gravidade não desapareceu:
astronauta e nave aceleram juntos ao redor da Terra.
\end{dica}

\begin{exercicios}
\nivel{1}{Conceitos}
\begin{questoes}
\item Um livro parado sobre a mesa recebe quais forças principais?
\item Calcule o peso de uma mochila de $6\un{kg}$, com $g=10\un{m/s^2}$.
\item Explique por que a força normal pode ser diferente do peso.
\item Se $\mu_e=\dec{0,40}$ e $N=50\un{N}$, qual é o atrito estático máximo?
\item Ao dobrar a distância entre duas massas, por qual fator muda a força gravitacional?
\nivel{2}{Aplicação}
\item Um bloco de $5\un{kg}$ desliza em piso horizontal com $\mu_c=\dec{0,20}$. Ache o atrito.
\item No item anterior, uma força horizontal de $20\un{N}$ puxa o bloco. Ache a aceleração.
\item A força entre duas massas é $36\un{N}$. Qual será o valor a uma distância três vezes maior?
\item Segundo a 2ª lei de Kepler, onde um planeta apresenta maior rapidez orbital?
\nivel{3}{Síntese}
\item Compare peso e massa e indique suas unidades.
\item Um planeta tem a mesma massa da Terra e raio duas vezes maior. Compare sua gravidade superficial com a terrestre.
\item Explique por que um satélite em órbita está em queda livre sem atingir imediatamente o solo.
\end{questoes}
\gabarito{2) $60\un{N}$; 4) $20\un{N}$; 5) $F/4$; 6) $10\un{N}$;
7) $2\un{m/s^2}$; 8) $4\un{N}$; 9) mais próximo do Sol; 11) $g/4$.}
\end{exercicios}


\end{document}
